\begin{table}
\centering
\caption{Robustness: FINBRA vs CAPAG}
\centering
\begin{tabular}[t]{lccc}
\toprule
  & M2: Composite & M4: Disaggregated & CAPAG grade\\
\midrule
Administrative capacity (composite) & -0.115*** &  & \\
 & (0.034) &  & \\
Education spending share &  & 1.565*** & \\
 &  & (0.182) & \\
Budget execution rate &  & 0.700** & \\
 &  & (0.343) & \\
Debt ratio &  & 0.027 & \\
 &  & (0.017) & \\
FPM dependence &  & -0.855*** & \\
 &  & (0.216) & \\
CAPAG grade (A=4 to D=1) &  &  & 0.002\\
 &  &  & (0.017)\\
\midrule
N & 5252 & 4742 & 3478\\
R2 & 0.523 & 0.539 & 0.551\\
R2 Adj. & 0.522 & 0.538 & 0.550\\
\bottomrule
\multicolumn{4}{l}{\rule{0pt}{1em}* p $<$ 0.1, ** p $<$ 0.05, *** p $<$ 0.01}\\
\multicolumn{4}{l}{\rule{0pt}{1em}Clustered SE by state in parentheses. Controls: log population, log GDP pc, log area, urban share, log distance to capital, biome FE, state fiscal autonomy. CAPAG model estimated on subsample with valid CAPAG rating (N = 3,116 vs 5,252 for FINBRA models). * p $<$ 0.1, ** p $<$ 0.05, *** p $<$ 0.01.}\\
\end{tabular}
\end{table}
